%% Diagramme d'états-transitions (stm) d'un volet roulant automatique.
%% Deux états, deux transitions déclenchées par la comparaison de N au seuil 614.
\documentclass[tikz,border=4pt]{standalone}
\usepackage{liaisons}
\begin{document}
\begin{tikzpicture}[x=1cm,y=1cm,
  cadre/.style={draw=black!70, line width=0.9pt},
  etat/.style={draw=solideB, line width=1.1pt, fill=white, rounded corners=5pt},
  fl/.style={-{Stealth[length=5pt]}, line width=1pt, solideA},
  ev/.style={font=\scriptsize, text=solideA, inner sep=2pt},
  leg/.style={font=\scriptsize, text=black!60, anchor=west, align=left, text width=1.9cm},
  amorce/.style={black!55, dash pattern=on 2pt off 2pt, line width=0.5pt}]

%% ================== cadre à onglet ====================================
\draw[cadre] (0,0) rectangle (5.3,4.4);
\draw[cadre] (0,3.95) -- (3.3,3.95) -- (3.55,4.18) -- (3.55,4.4);
\node[font=\scriptsize, anchor=west, inner sep=3pt] at (0.1,4.17)
     {\texttt{stm} [Machine à État] Volet};

%% ================== pseudo-état de démarrage ==========================
\fill[black] (0.75,3.65) circle (0.11);
\draw[fl] (0.86,3.6) -- (1.7,3.27);

%% ================== état « Volet ouvert » =============================
\draw[etat] (1.4,2.55) rectangle (3.8,3.25);
\node[font=\small] at (2.6,2.9) {Volet ouvert};

%% ================== état « Volet fermé » (avec activités internes) ====
\draw[etat] (0.55,0.5) rectangle (4.65,1.78);
\node[font=\small] at (2.6,1.55) {Volet fermé};
\draw[solideB, line width=0.8pt] (0.55,1.34) -- (4.65,1.34);
\node[font=\scriptsize, anchor=west] at (0.7,1.13) {\texttt{entry} / Éclairer voyant rouge};
\node[font=\scriptsize, anchor=west] at (0.7,0.76) {\texttt{exit} / Éteindre voyant rouge};

%% ================== transitions =======================================
\draw[fl] (1.9,2.55) -- (1.9,1.81);
\node[ev, anchor=east] at (1.82,2.25) {$N > 614$};
\draw[fl] (3.5,1.78) -- (3.5,2.52);
\node[ev, anchor=west] at (3.58,2.25) {$N \le 614$};

%% ================== légende ===========================================
\draw[amorce] (0.9,3.65) -- (5.55,3.65);
\node[leg] at (5.6,3.65) {Pseudo état de démarrage};
\draw[amorce] (3.85,2.9) -- (5.55,2.9);
\node[leg] at (5.6,2.9) {État};
\draw[amorce] (3.52,1.95) -- (5.55,1.95);
\node[leg] at (5.6,1.95) {Transition (événement)};
\end{tikzpicture}
\end{document}
